\section{The original interacting clock on coherent routes}
% Chapter VI: original interacting source, conditional proper-time interface.
\subsection{The full interacting source as a route marker}
\label{sec:ch6-clock-routes}
A route-dependent clock requires a receiving rule connecting the local
interacting dynamics to spacetime transport. The following rule is conditional~\cite{Zych2011}:
the same isolated, co-moving source--field Hamiltonian acts along each prescribed
timelike route, with its Chapter I time chart replaced by accumulated proper
time. External beam splitters, recombination, routing forces, and the spacetime
geometry are controls; they are not derived from the source Hamiltonian.
Tidal changes of the internal Hamiltonian, radiation caused by the controls,
and inaccessible environmental records require additional terms or a larger
marker system. A field shared nonlocally by both arms does not automatically
admit this isolated co-moving reduction.

Retain the complete source space $\mathbb C^4\otimes\mathbb C^4$ and the two
polarizations of the continuum Fock field. In the Chapter I dimensionless
chart the rest generator is
\begin{equation}
\begin{split}
 H_0={}&\epsilon(S+Q_L+Q_R+Q_LQ_R)+H_f\\
 &-\epsilon(D\phi_++C\phi_\times)\\
 &+\epsilon^2\mu(D^2+C^2),\\
 D={}&Q_L-Q_R,\qquad C=iSD,\qquad \mu=(4R)^{-1}.
\end{split}
\label{eq:ch6-full-clock-generator}
\end{equation}
This is the original interacting form-sum Hamiltonian, not a bare two-level
substitution. Its six two-state lanes and four unforced diagonal states are
all retained. In a chosen prepared lane $a<b$, $d=q_a-q_b<0$,
$g=\epsilon d$, and the restriction is
$H_0=\epsilon eI+\epsilon S+H_f-g(\sigma_z\phi_++\sigma_y\phi_\times)+cI$,
where $e=q_a+q_b+q_aq_b$ and $c=2g^2\mu$. Both the even energy and the two
counterterms remain in the generator.

For route $j$, define
\begin{equation}
 \begin{gathered}
 \eta_j=s_0\int_{\gamma_j}d\tau_{\rm proper},\qquad
 U_j=e^{-iH_0\eta_j},\\
 U_{\rm ctrl}=\sum_{j=1}^2|j\rangle\langle j|\otimes U_j.
 \end{gathered}
 \label{eq:ch6-controlled-clock}
\end{equation}
Here $s_0=1/t_u$ in the declared Chapter I chart; no new SI calibration is
inferred. With $I_*=\hbar$, the physical rest energy
is $(\hbar/t_u)H_0$. Conditional on this transport law, the actual clock
coordinate obeys $d\theta/dt=\nu(\eta_j)(d\tau_{\rm proper}/dt)/t_u$ wherever
its certified inverse exists. This does not extend the Chapter I finite-window
rate certificate beyond its stated preparation and interval.

Let the route input be $\sqrt{p_1}|1\rangle+e^{i\alpha}\sqrt{p_2}|2\rangle$,
$p_1+p_2=1$, initially uncorrelated with a marker state $\rho$. Tracing the
marker gives route coherence proportional to
\begin{equation}
\begin{gathered}
\Gamma(\delta\eta)=\operatorname{Tr}(\rho U_2^\dagger U_1)
 =\operatorname{Tr}(\rho e^{-iH_0\delta\eta}),\\
 \delta\eta=\eta_1-\eta_2.
\end{gathered}
\label{eq:ch6-route-overlap}
\end{equation}
An ideal variable-phase recombiner measures fringe visibility
$\mathcal V=2\sqrt{p_1p_2}|\Gamma|$. The standard Born measurement rule is assumed in this calculation, not
derived from the original Write histories. The phase of $\Gamma$ shifts the fringes;
its modulus controls their contrast. Extra uncontrolled route records multiply
or otherwise modify this coherence and must be included explicitly.

For balanced priors, the operational which-route distinguishability in the
full marker is
\begin{equation}
\begin{gathered}
\mathcal D=\tfrac12\|U_1\rho U_1^\dagger-U_2\rho U_2^\dagger\|_1,
 \\
 \mathcal D^2+|\Gamma|^2\leq1.
\end{gathered}
\label{eq:ch6-duality}
\end{equation}
Equality holds for a pure initial full marker with unrestricted joint
measurements. For a mixed marker the inequality follows by purifying $\rho$:
the two purified states have overlap $\Gamma$, and partial trace contracts
trace distance. Restricted access to the source alone can only decrease
$\mathcal D$. For unequal priors the corresponding statement uses
$\|p_1\rho_1-p_2\rho_2\|_1$ and $\mathcal V$.
In particular, if $[\rho,H_0]=0$, then $\rho_1=\rho_2$ and $\mathcal D=0$,
although $|\Gamma|$ may be small. Mixed-state phase averaging is not itself
accessible which-route information in that marker. Nor does the phase-clock
expectation $z=\langle Z\rangle$ determine $\Gamma$: one is an expectation
of an evolved observable, the other a full-state propagator overlap.

\subsection{The needed spectral operand and its domain}
\label{sec:ch6-spectral-clock-operand}
Let $P_{H_0}$ be the spectral measure of the full interacting Hamiltonian.
The new original-source operand needed for exact route contrast is
\begin{equation}
\begin{gathered}
d\mathcal P_\rho(E)=\operatorname{Tr}[\rho P_{H_0}(dE)],\\
 \Gamma(s)=\int e^{-iEs}\,d\mathcal P_\rho(E).
\end{gathered}
\label{eq:ch6-spectral-operand}
\end{equation}
Chapter I's bounds on $M$, $\Omega$, $z$, and $\nu$ do not provide this
entire measure or the exact Loschmidt amplitude. A new spectral enclosure,
a controlled evaluation of this matrix element, or an overlap enclosure
with a stated error is needed for a numerical full-continuum fringe curve.
The matched hard mean energy $\bar E$ is also required for its absolute
reference phase; a gap enclosure alone does not determine that phase.

If $\int E^2d\mathcal P_\rho(E)<\infty$, set
$\sigma_H^2=\int(E-\bar E)^2d\mathcal P_\rho(E)$. Then
\begin{equation}
\begin{gathered}
|\Gamma(s)|^2=1-\sigma_H^2s^2+o(s^2),\\
 |\Gamma(s)|=1-\tfrac12\sigma_H^2s^2+o(s^2).
\end{gathered}
\label{eq:ch6-short-time-contrast}
\end{equation}
These coefficients use the full prepared joint-energy distribution.
Neither the bare exchange gap $2\epsilon$ nor the calibration expectation
$E_{\rm ref}=573049$ is an energy variance. The latter belongs to a different
placement-basis calibration preparation. For a pure state the requisite
condition is $\Psi\in D(H_0)$, not merely the energy form domain
$D(H_f^{1/2})$. Chapter I explicitly has ultraviolet-divergent
$\int_0^\infty|u_R|^2d\omega$; consequently a generic bare vacuum argument
cannot supply this second moment.

When this finite-variance condition holds, the standard quantum speed
bound~\cite{Fleming1973,DeffnerCampbell2017} follows because the projective Hilbert-space
speed of a purification evolving under $H_0$ is $\sigma_H$. Thus
\begin{equation}
\begin{aligned}
 \arccos|\Gamma(s)|&\leq\sigma_H|s|,\\
 |\Gamma(s)|&\geq\cos(\sigma_H|s|),
 \quad \sigma_H|s|\leq\pi/2.
\end{aligned}
 \label{eq:ch6-speed-bound}
\end{equation}
Any zero of the amplitude requires $|s|\geq\pi/(2\sigma_H)$.
For a mixed initial marker this is a bound on the specified amplitude (or its
purification), not a claim that the reduced marker becomes distinguishable
at that time. The full pure-state marker also obeys
$\mathcal D\leq\sin(\sigma_H|s|)$ on this interval.

\paragraph{A finite full-H second moment for the actual declared preparation.}
The ultraviolet qualification does not invalidate the particular hard-branch
preparation that Chapter I actually declares. Let $H_{\rm ref}$ be its
matched hard-interacting Hamiltonian, with the same even scalar and soft
counterterm as $H_0$, and let
\begin{equation}
\begin{aligned}
 \psi_k&=\Phi_k\otimes|0_s\rangle,
 &\Psi_0&=(\psi_0+\psi_1)/\sqrt2,\\
 H_{\rm ref}\psi_k&=E_k\psi_k,
 &E_1-E_0&=\Omega.
\end{aligned}
\end{equation}
The normalized exact hard eigenvectors are operator-domain vectors of their
hard Hamiltonian. The reference preparation is therefore in
$D(H_{\rm ref})$. Write
\begin{equation}
\begin{aligned}
 V_s&=H_0-H_{\rm ref}\\
    &=-g(\sigma_z\phi_{+,s}+\sigma_y\phi_{\times,s}),\\
 v_s&=\int_0^\lambda|u_R(\omega)|^2\dd\omega
 \leq v_s^{\rm up}=\frac{R^4\lambda^6}{675\kappa\pi}.
\end{aligned}
\end{equation}
On this soft vacuum, $V_s\Psi_0$ is a Hilbert-space vector with exactly one
soft photon. The two polarization insertions are orthogonal and both Pauli
operators preserve norm, giving $\|V_s\Psi_0\|^2=2g^2v_s$.
The form identity against every test vector in the common form domain is
represented by the Hilbert vector $H_{\rm ref}\Psi_0+V_s\Psi_0$.
The representation theorem for the closed form therefore proves
$\Psi_0\in D(H_0)$ and $H_0\Psi_0=H_{\rm ref}\Psi_0+V_s\Psi_0$.
Orthogonality of zero and one soft photon sectors now gives the new exact
identity
\begin{equation}
\begin{gathered}
\bar E=\tfrac12(E_0+E_1),\\
 \boxed{\ \sigma_H^2=\frac{\Omega^2}{4}+2g^2v_s\ }
 \leq\frac{\Omega^2}{4}+2g^2v_s^{\rm up}.
\end{gathered}
\label{eq:ch6-prepared-full-variance}
\end{equation}
This is a full interacting variance, despite its expression in matched
reference data. It relies on the exact hard eigenbranch preparation; it is
not a general domain assertion about arbitrary coherent source states.
It also does not assert $\Psi_0\in D(H_0^2)$ or existence of higher moments.
The signed $g$ is retained in the Hamiltonian; its square appears here because
this is a norm. A rigorous numerical speed enclosure follows by using the
original certificate's upper bounds for $\Omega$, $|g|$, $R$, and $\lambda$,
with their declared dependencies, rather than rounded lane numbers.

\paragraph{Six-lane consequences of the retained inputs.}
Table~\ref{tab:ch6-route-clock-bounds} uses the original rational enclosures
for $\Omega$, signed $g$, $v_s^{\rm up}$ and $T_0=2\pi/\Delta$.
The new rational operations are recorded in the companion check output;
each printed bound is verified outward by exact rational comparison.
These are independently derived enclosures, not a new native-certificate run.
The lower orthogonality ratio is slightly above one half in every lane,
while the reference-amplitude error at $T_0$ is at most $4.54\times10^{-11}$.
This does not imply that an exact zero or a perfect revival occurs.
\begin{table}[ht]
\centering\scriptsize
\begin{tabular}{c c c c}
\hline
Lane & $10^6\sigma_H$ upper & error at $T_0$ upper & $s_\perp/T_0$ lower\\
\hline
01 & 1.745051466079 & $1.759\times10^{-12}$ & 0.500000001207\\
02 & 1.745051452191 & $3.644\times10^{-12}$ & 0.500000005186\\
03 & 1.745048665211 & $4.536\times10^{-11}$ & 0.500000803726\\
12 & 1.745051465446 & $1.886\times10^{-12}$ & 0.500000001389\\
13 & 1.745048878445 & $4.360\times10^{-11}$ & 0.500000742629\\
23 & 1.745049097788 & $4.172\times10^{-11}$ & 0.500000679781\\
\hline
\end{tabular}
\caption{Full interacting hard-branch preparation: variance-speed and
matched-reference bounds. The error is an amplitude error, and $s_\perp$
is any first possible overlap zero, in the dimensionless proper-time chart.}
\label{tab:ch6-route-clock-bounds}
\end{table}

\paragraph{An independent overlap enclosure from the same preparation.}
The reference overlap is
$\Gamma_{\rm ref}(s)=e^{-i\bar Es}\cos(\Omega s/2)$.
The reference orbit is a finite sum of operator-domain eigenvectors, and its soft insertion is Hilbert-space continuous. Strong Hilbert-space Duhamel, with the full interacting propagator left intact, gives
\begin{equation}
\begin{aligned}
 &\|(e^{-isH_0}-e^{-isH_{\rm ref}})\Psi_0\|\\
 &\quad\leq\int_0^{|s|}\|V_s e^{-iuH_{\rm ref}}\Psi_0\|\dd u
 =|s|\sqrt{2g^2v_s}.
\end{aligned}
\end{equation}
The same orthogonal soft-sector norm is constant for every reference time;
negative $s$ is handled by reversing the integral. Consequently
\begin{equation}
\begin{aligned}
 \varepsilon_s&=|s|\sqrt{2g^2v_s^{\rm up}},\\
 |\Gamma(s)-e^{-i\bar Es}\cos(\Omega s/2)|&\leq\varepsilon_s,\\
 |\Gamma(s)|&\geq
 \max\{0,|\cos(\Omega s/2)|-\varepsilon_s\}.
\end{aligned}
 \label{eq:ch6-overlap-enclosure}
\end{equation}
This is a newly derived overlap estimate using Chapter I's exact preparation,
not an inference from its phase-rate certificate. It preserves all photon
sectors in the actual propagation. The bound can become uninformative at
long times; it does not certify phase-clock invertibility there.

For another preparation lacking a verified second moment, a safe alternative
is spectral truncation. For a bounded interval $I$, let
$p=\operatorname{Tr}[\rho P_{H_0}(I)]>0$ and normalize its restricted spectral
measure. Its amplitude $\Gamma_I$ satisfies
$|\Gamma-p\Gamma_I|\leq1-p$ and hence
$|\Gamma|\geq\max\{0,p|\Gamma_I|-(1-p)\}$.
A finite-variance speed bound for the restricted measure may then be used,
but the tail probability $1-p$ must itself be established for the original
state. A finite matrix cutoff without that tail is a comparator, not a
continuum certificate.

\subsection{Action phases, energy offsets, and route records}
\label{sec:ch6-route-phase-account}
The coherent amplitude of a controlled route may contain a scalar factor
$e^{iS_{{\rm ext},j}/\hbar}U_j$. Its scalar interference phase is
$(S_{{\rm ext},1}-S_{{\rm ext},2})/\hbar$, combined with the phase of
$\Gamma$. This accounting must partition the generator once. If an assigned
rest mass $m_0$ is excluded from $H_0$, its point-particle contribution is
$S_{{\rm mass},j}=-m_0c^2\int_{\gamma_j}d\tau_{\rm proper}$.
If that same energy is already included in the internal generator, adding
this phase again double counts it. The same warning applies to an even
source-energy offset and the scalar counterterm in a fixed lane.

Explicitly, replacing $H_0$ by $H_0+bI$ changes
$\Gamma(s)$ to $e^{-ibs}\Gamma(s)$ without changing visibility, the
same-history expectation $z$, or which-route distinguishability. Thus an
energy offset invisible to the original clock-rate equation can produce a
relative route phase $-b(\eta_1-\eta_2)$. An unchanged total amplitude requires the
compensating assignment
$S_{{\rm ext},j}\mapsto S_{{\rm ext},j}+\hbar b\eta_j$.
A lane-dependent even operator is not a common scalar on a superposition of
lanes and can affect its full spectral distribution. The original signed
Write integral, the quantum rest-generator phase, and a relativistic
point-particle action are therefore distinct operands until a receiving
identification and a non-overlapping energy account are supplied.

Finally, distinct ordered route or Write records need not imply distinct
accessible quantum states. If two histories induce the same observable
state and leave no coupled record, an extra history label alone changes no
measurement probability. If a physical record is produced, it belongs in
the marker or environment and its overlap controls the residual coherence.
This allows different histories with equal endpoint and equal signed action
without silently creating an additional observable or a new clock law.

\paragraph{Independent finite-source comparator.}
The companion executable \texttt{check\_clock\_routes.py} retains the exact
$16$-state source matrices, signed $D,C$, the full even source energy, both
polarizations, and both discrete counterterms. It checks source calibration,
controlled-route coherence, pure and mixed duality, the short-time
coefficient, the hard/soft variance identity, the Duhamel enclosure, and
scalar-phase accounting. Its finite oscillator modes are explicitly assigned
comparison data. Their tests neither replace the continuum smearing nor
certify its spectral measure; the continuum results above follow from the
stated form-domain and hard-branch arguments.

